% !TEX root = ../Grunddatei.tex
\section{Related works}

When looking at the domain of remote sensing, the inclusion of artificial intelligence is advancing rapidly. When searching for “satellite image” and “deep learning” on the research-sharing platform \textit{arXiv.org}, the number of published articles exceeds one hundred in 2025 alone. The developments in image analysis using artificial intelligence can improve existing use cases as well as introduce new ones. A small section of related articles will be presented in this chapter and can be seen as representative for other related works.

At the 17th Python in Science Conference in 2018, Virginia Ng and Daniel Hofmann presented a methodology to implement scalable tools for extracting features in remote sensing images, using deep learning methods. The proposed approach can be used for a variety of tasks, for example object detection, segmentation and feature mapping. It also covers various features, ranging from roads and bikes to bodies of water or even glacies. The paper includes an overview of the basic pipeline, that includes the data preparation as well as implementation of different deep learning models such as U-Net or YOLO. The authors also presented specific adaptations of the methodology for the aerial and satellite imagery domain. Last, they discussed further efforts to improve the proposed deep learning models, so that new applications can also be covered \cite{relWork1}.

The second paper, written by Lei He et al, propose an improved deep learning approach for accurate road segmentation from remote sensing imagery. Due to the improved segmentation precision and reduced computational complexity, the new technology can be used to extract roads from satellite imagery. The authors introduce a modified architecture called D-DenseNet, which includes parts of both D-LinkNet and DenseNet. The proposed approach demonstrates its suitability for precise and computationally efficient road extraction \cite{relWork2}.

Chenao Zhao et al challenge current models when faced with vehicle detection tasks in satellite imagery, by introducing a novel vehicle detection called SatDetX-YOLO. The new approach is based on YOLOv8 and is characterized by several architectural changes that stem from combinations with other models. For instance, the model includes a backbone network with FasterNet as well as a Deformable Attention Module. When comparing the results between different models, a slight increase in accuracy, precision, recall and mean average precision can be observed (around 3 percent respectively) \cite{relWork3}.

The last presented paper, by Andrew Campbell, Alan Both and Qian Sun, explores the possibilities of automated classification tasks in Google Street View images using machine learning technologies. Focus of their work lies on the detection of traffic signs. The authors developed a custom object detection model to classify certain street signs from images captured at intersections and get their approximate locations. The implementation shows notable detection and classification accuracies of around 96 and 98 percent respectively as well as the ability to be easily scalable. Due to those characteristics and the workflow using open-source data, the proposed solution shows great potential for real-world applications \cite{relWork4}.

From what we can infer, our approach is relatively novel, as we could not find articles that follow a similar approach for one-way street identification. Although, a lot of research has been done for satellite images and street view panoramas respectively, combined solutions using both seem rather unique. The same appears to be true for segmentations of road intersections in satellite imagery, which seems to be a less explored research topic.